%@ pages 125-135
%@ status translated
% Original section 12, "La famille universelle", pp. 125--135.
% Section 12, "The universal family", original pp. 125--135.
%
% Kept as printed, recorded here for the next reader:
%  * after item (ii) the authors write "si (b) est vrai"; the list above is
%    numbered (i), (ii), so (b) is the printed form of (ii);
%  * "$B|_{z/\Zz}$" (four times, original p. 132) where $B$ is a functor on
%    $\Yy$, so that $z/\Yy$ is meant; printed as is;
%  * the prose says that a morphism between matching objects "est fibrant"
%    where "est une fibration" is meant (original p. 134);
%  * in the recall paragraph, "fibrant over $\Xx$" while the site has just
%    been renamed $\Yy$; printed as is.
% One clear typographical slip has been repaired, with a comment at the spot:
% the stray "\ldots" written as plain "ldots" (original p. 131).

\clearpage
\section{The universal family}
\label{universellepage}

For the moment let $\Xx$ be a category equipped with the coarse topology.
One has a morphism of $n+1$-prestacks
$$
\Psi : nSeCHAMP(\Xx ) \times \Xx ^o \rightarrow nSeCAT.
$$
This morphism may be seen as the \emph{universal family} of $n$-categories
over $nSeCHAMP(\Xx ) \times \Xx$. We define $\Psi$ by setting
$$
\Psi (A,X):= A(X);
$$
and for $f:X\rightarrow Y$ in $\Xx $ we use the natural morphism
$$
\Gamma \underline{Hom}(A,B) \times \{ f\} \rightarrow \underline{Hom}(A(Y),
B(X))
$$
to define the action of $\Psi$ on morphisms. On the left, this is the
internal $\underline{Hom}$ of Segal $n$-prestacks; and, on the right, it is
the internal $\underline{Hom}$ of Segal $n$-precategories.

Let $nSeCAT'$ be the fibrant replacement of $nSeCAT$. By composition one
obtains a morphism
$$
nSeCHAMP(\Xx ) \times \Xx ^o\rightarrow nSeCAT',
$$
which corresponds, by definition of the internal $\underline{Hom}$ (of Segal
$n+1$-prestacks), to a morphism
$$
\Phi : nSeCHAMP(\Xx ) \rightarrow \underline{Hom}(\Xx ^o, nSeCAT').
$$

The following result makes the link between these definitions and the notions
of stack proposed in \cite{limits}.


\begin{theorem}
\label{correlation}
Let $nSeCAT'$ be a fibrant replacement of $nSeCAT$.
If the topology of $\Xx$ is
coarse then the morphism
$\Phi$ induces an equivalence between
$nSeCHAMP(\Xx )$ and the
Segal $n+1$-category $\underline{Hom}(\Xx ^o, nSeCAT')$.

If $\Xx$ is equipped with another topology
$\Gg$ then $\Phi$ induces an equivalence between \newline
$nSeCHAMP(\Xx )$ and the full subcategory
of
$\underline{Hom}(\Xx ^o, nSeCAT')$ consisting of the morphisms
$F:\Xx ^o\rightarrow nSeCAT'$ which correspond to Segal $n$-stacks.

% modified in v3 (the preceding and the following sentence)
For the case of non-Segal $n$-stacks this subcategory may be
identified as that of the morphisms which
satisfy the descent condition
of \cite{limits} 6.3, namely that the arrow
$$
\lim _{\leftarrow} F|_{\Xx /X} \rightarrow
\lim _{\leftarrow} F|_{\Bb}
$$
be an equivalence for every sieve $\Bb \subset \Xx /X$ of $\Gg$.
\end{theorem}

In this section (Proposition \ref{pleinfidele} below) we shall
only
prove that $\Phi$ is fully faithful. The essential surjectivity of
$\Phi$ in the first case of the theorem
will be Theorem \ref{strictif3}
further down. The second case will then follow
(once Theorem \ref{strictif3} is granted) from Corollary
\ref{critereaveclim}.

We note first that if $\Gg$ denotes the
topology of $\Xx$,
then the $\Gg$-fibrant $n$-prestacks are in particular fibrant for the
coarse topology, whereas the definition of the morphisms in
$nSeCHAMP(\Xx ^{\Gg})$
does not involve $\Gg$. Hence $nSeCHAMP(\Xx ^{\Gg} )$ is a full
subcategory
of $nSeCHAMP(\Xx ^{\rm gro})$, where $\Xx ^{\rm gro}$ is the category
$\Xx$ with
the coarse topology. Thus the property that $\Phi$ be fully
faithful reduces to the case of the coarse topology
and we may ignore the topology $\Gg$
for the rest of this
section.

To emphasise the fact that the topology is not being considered, we change
the notation for $\Xx$ to $\Yy$: we therefore consider a category $\Yy$ with
its coarse topology.

Recall that the objects of $nSeCHAMP(\Yy )$ are the Segal $n$-prestacks
fibrant over $\Xx$, and that if $A$ and $B$ are two such objects, one has,
$$
nSeCHAMP (\Yy )_{1/}(A,B) := \Gamma \underline{Hom}(A,B).
$$
This fibrant Segal $n$-precategory represents the functor which to a
Segal $n$-precategory $E$ associates the set $Hom (A\times \underline{E}, B)$.
On the other hand, if $A$ is fibrant then the values $A(X)$ are fibrant, i.e.
objects of
$nSeCAT$. The fibrant Segal $n$-prestack $A$ therefore induces a morphism
$\Yy ^o \rightarrow  nSeCAT$ which one may compose with the fibrant
replacement
to obtain a morphism $\Yy ^o \rightarrow  nSeCAT'$.
If $E$ is a Segal $n$-precategory $E$, to give a
morphism
$$
E\rightarrow \underline{Hom}(\Yy ^o, nSeCAT')_{1/}(A,B)
$$
amounts to giving a morphism
$$
F:\Yy ^o\times \Upsilon (E) \rightarrow nSeCAT'
$$
with
$$
F|_{\Yy ^o \times \{ 0\}} = A,\;\;\;\;
F|_{\Yy ^o \times \{ 1\} } = B.
$$
(See \cite{limits} for the notation $\Upsilon$.)
Now, the
data of such a
strict morphism (i.e. with values in $nSeCAT$ instead of $nSeCAT'$)
$$
F: \Yy ^o\times \Upsilon (E) \rightarrow nSeCAT
$$
with
$$
F|_{\Yy ^o \times \{ 0\}} = A,\;\;\;\;
F|_{\Yy ^o \times \{ 1\}} = B
$$
amounts to the data, for every $y\in \Yy$, of a morphism
$E\times A(y)\rightarrow B(y)$ such that for $y\rightarrow z$ in $\Yy$, these
morphisms are strictly compatible with the restrictions in $A$ and $B$.
In other words, one obtains the equality
$$
\underline{Hom}(\Yy ^o, nSeCAT)_{1/}(A,B) = \Gamma \underline{Hom}(A,B).
$$
Consequently the inclusion $nSeCAT \rightarrow nSeCAT'$ induces a morphism
$$
\Gamma \underline{Hom}(A,B) \rightarrow
\underline{Hom}(\Yy ^o, nSeCAT')_{1/}(A,B),
$$
and the full faithfulness of the functor
$\Phi$ is equivalent to the following proposition.

\begin{proposition}
\label{pleinfidele}
Let $A$ and $B$ be two fibrant Segal $n$-prestacks on $\Yy$. Then the
morphism
$$
\Gamma \underline{Hom}(A,B) \rightarrow
\underline{Hom}(\Yy ^o, nSeCAT')_{1/}(A,B).
$$
is an equivalence of Segal $n$-categories.
\end{proposition}
{\em Proof:} We first introduce some notation.
Let $Y$ be a Segal $n+1$-category and $E$ a Segal $n$-precategory. We
define the Segal $n+1$-precategory $Y^{\otimes E}$ which has the same objects
as $Y$ by setting (for $p\geq 1$)
$$
Y^{\otimes E}_{p/}(y_0,\ldots , y_p):=
Y_{p/}(y_0,\ldots , y_p) \times E.
$$
We are interested in this construction mainly when $Y$ is
our base $1$-category $\Yy$, or else for $Y=\Yy \times I$.


If $\Zz \subset \Yy$ is an inclusion of $1$-categories having the same
objects, if $C'$ is a fibrant $n+1$-category, and if $A : \Zz \rightarrow
C'$ is a
morphism, we introduce the fibrant Segal $n$-category $Diag (\Yy , \Zz ,
A; C')$, whose objects are the extensions of $A$ to $\Yy$, and which
is characterised by the formula
$$
Hom _{nSePC}(E,Diag (\Yy , \Zz , A;C'))=\{ f: \Yy ^{\otimes E} \rightarrow C',
\;\;\; f|_{\Zz ^{\otimes E}} = A|_{\Zz ^{\otimes E}} \} .
$$
Indeed, one verifies that the right-hand side defines a functor
in $E$ which
transforms colimits into limits. This functor is therefore
representable by a Segal $n$-precategory which we denote
$Diag (\Yy , \Zz ,
A;C')$.
Since $C'$ is
fibrant, our functor has the extension property for
trivial cofibrations $E\subset E'$, and hence $Diag (\Yy , \Zz , A;C')$ is
fibrant.

Let $\nu (\Zz )\subset \nu (\Yy )$ and $\nu (C')$ be the nerves of these Segal
$1$- or
$n+1$-categories, regarded as Segal $n$-prestacks
over $\Delta$. We note that $\nu (C' )$ is fibrant over
$\Delta$. One may verify
that the Segal $n$-category $Diag (\Yy , \Zz , A;C')$  is the fibre of the
morphism
$$
\underline{Hom}(\nu (\Yy ), \nu (C' ))
\rightarrow
\underline{Hom}(\nu (\Zz ), \nu (C' ))
$$
over $A$. What is involved here are the internal $\underline{Hom}$'s of Segal
$n$-prestacks over $\Delta$.

Let now $A,B: \Yy \rightarrow C'$ and let us denote by $A\sqcup B$ the
corresponding morphism
$$
A\sqcup B:\Yy \times \{ 0,1\} = \Yy \sqcup \Yy \rightarrow C'.
$$
One has the formula
$$
\underline{Hom}(\Yy , C')_{1/}(A,B) = Diag (\Yy \times I, \Yy \times \{ 0,1\} ,
A\sqcup B ; C').
$$
To prove full faithfulness, we shall apply this formula
with
$C'=nSeCAT'$.

We shall first modify things slightly.
In the situation of the preceding
paragraph, let $\Uu$ be the category whose objects are the pairs of
sequences $(y_0,\ldots , y_p; z_0, \ldots , z_q)$ equipped with morphisms
$y_i\rightarrow y_{i+1}$,
$z_i\rightarrow z_{i+1}$ and also $y_p\rightarrow z_0$; and where the
morphisms
are the inclusions and the degeneracies of sequences.
In other words, $\Uu$ is the \emph{category of the simplices of $\Yy \times I$}
which lie neither in $\Yy \times \{ 0\}$ nor in $\Yy \times
\{ 1\}$. Let us remark that one has a forgetful functor
$$
(y_0,\ldots , y_p; z_0, \ldots , z_q)\mapsto
[\{ 0,\ldots , p\} , \{ 0,\ldots ,
q\} ]
$$
from
$\Uu$  to $\Delta \times \Delta$.
One may define
an $n$-prestack $F = F(A,B; C')$ on
$\Uu$ by taking for
$F(y_0,\ldots , y_p,z_0,\ldots ,z_q)$
the fibre over
$(f,g)$ of the morphism
$$
C'_{p+q+1/}(A(y_0),\ldots , A(y_p),B(z_0),\ldots , B(z_q))
\rightarrow
$$
$$
C'_{p/}(A(y_0),\ldots , A(y_p))\times
C'_{q/}(B(z_0),\ldots , B(z_q))
$$
where $f$ (resp. $g$) is the image of the element $(y_0,\ldots , y_p)$
(resp. $(z_0,\ldots , z_q)$) of the nerve of $\Yy$ under the morphism $A$ (resp. $B$).


Let us observe here the following property:

$(\ast ):$ for every $(y_0,\ldots , z_q)$,
the restriction morphism
$$
F(y_0,\ldots , y_p,z_0,\ldots ,z_q)\rightarrow F(y_p,z_0)
= C'_{1/}(A(y_p), B(z_0))
$$
corresponding to the arrow
$$
(y_p; z_0)\rightarrow (y_0,\ldots ,y_p; z_0,\ldots , z_q)
$$
% The original prints "z_0,ldots , z_q" here; the backslash is missing.
of $\Uu$, is
an equivalence of Segal $n$-categories.

To see this, one must use
the hypothesis that $C'$ is fibrant, which makes the fibre product in the
definition of $F$ also a homotopy fibre product.

We shall now finish the proof in progress by assuming
that $F$ is fibrant over $\Uu$, and we shall devote the end of
the present section to this problem.
Let us observe now that one has
$$
Diag (\Yy \times I, \Yy \times \{ 0,1\} ,
A\sqcup B ; C') = \Gamma (\Uu , F).
$$
Indeed, if $E$ is a
Segal $n$-precategory, the data of a morphism
$E\rightarrow \Gamma (\Uu ,F)$
is equivalent to the data, for every object
$(y_0,\ldots , y_p,z_0,\ldots ,z_q)$ of $\Uu$, of a morphism
$$
E\rightarrow
C'_{p+q+1/}(A(y_0),\ldots , A(y_p),B(z_0),\ldots , B(z_q))
$$
compatible with the restrictions and
degeneracies of sequences,
and compatible with the transition morphisms of $A$ and $B$
over the sequences $(y_0,\ldots , y_p)$ and $(z_0,\ldots , z_q)$. Such
data is again equivalent to that of a
morphism $(\Yy \times I)^{\otimes E}\rightarrow C'$
compatible with $A\sqcup B$, whence the formula.

Let $Arr(\Yy )$ be the category whose objects are the arrows
$x\rightarrow y$ of $\Yy$ and whose morphisms from
$x\rightarrow y$ to $x'\rightarrow y'$ are the commutative diagrams
$$
\begin{array}{ccc}
x & \rightarrow & x' \\
\downarrow && \downarrow \\
y& \leftarrow & y' .
\end{array}
$$
One has a functor $\varphi : \Uu \rightarrow Arr (\Yy )$ defined by
$$
\varphi : (y_0,\ldots , y_p,z_0,\ldots , z_q)\mapsto (y_p\rightarrow z_0).
$$

We fix from now on $C'=nSeCAT'$, and we suppress the reference to
$C'$ in
$Diag$. Moreover we set $F:=F(A,B; nSeCAT')$,
where $A$ and $B$ are from now on
strict functors $\Yy \rightarrow nSeCAT$,
i.e. Segal $n$-prestacks on $\Yy ^o$ with
$B$
fibrant as an $n$-prestack on $\Yy^o$.

Let $\overline{hom}(A,B)$ be the $n$-prestack on $Arr (\Yy )$ which to
$y\rightarrow z$ associates the $n$-category \linebreak
$\underline{Hom}(A(y),B(z))$ (given our definition of the
morphisms of $Arr(\Yy )$ this
is indeed a contravariant functor on $Arr(\Yy )$). We have chosen
the notation $\overline{hom}$ to suggest
that this is a kind of ``external $Hom$''.

If we had defined $F$ with $nSeCAT$ instead of $nSeCAT'$
we would have had
$$
F(A,B, nSeCAT) = \varphi ^{\ast} (\overline{hom}(A,B)).
$$
Instead of that, with
$F:= F(A,B; nSeCAT')$
and the inclusion
$nSeCAT\rightarrow nSeCAT'$ one has a morphism
$$
\varphi ^{\ast} (\overline{hom}(A,B))\rightarrow F,
$$
which is, object by object over $\Uu$, an equivalence of
Segal $n$-categories (thanks to the equivalence $(\ast )$
seen above). By adjunction one obtains a morphism
$$
\overline{hom}(A,B)\rightarrow \varphi _{\ast}(F).
$$
On the other hand one has the formulas
$$
\underline{Hom}(\Yy , nSeCAT')_{1/}(A,B)=\Gamma (\Uu ,F) = \Gamma (Arr(\Yy ),
\varphi _{\ast} (F))
$$
and
$$
\Gamma \underline{Hom}(A,B)= \Gamma (   Arr(\Yy ), \overline{hom}(A,B)),
$$
which induce a commutative diagram
$$
\begin{array}{ccc}
\Gamma \underline{Hom}(A,B)&= &\Gamma ( Arr(\Yy ), \overline{hom}(A,B))\\
\downarrow && \downarrow \\
\underline{Hom}(\Yy , nSeCAT')_{1/}(A,B)&=& \Gamma (Arr(\Yy ),
\varphi _{\ast} (F)).
\end{array}
$$
To obtain the proposition,
i.e. that the left-hand vertical arrow is an equivalence, it suffices to
show that the right-hand vertical arrow is one. For this, it would
suffice  (given the fact that $\varphi _{\ast} (F)$ is fibrant
on $Arr(\Yy )$
cf \S 4)
to show the following two things:
\newline
(i)\,\, the morphism $\overline{hom}(A,B)\rightarrow \varphi _{\ast}(F)$ is
an equivalence object by object over $Arr(\Yy )$; and
\newline
(ii)\,\, $\overline{hom}(A,B)$ is fibrant on $Arr(\Yy )$.

In fact we do not know whether (b) is true, but we shall show a
weaker variant:
\newline
(ii)'\,\, for every $A$ there exists an equivalence $A'\rightarrow A$ such that
$\overline{hom}(A',B)$ is fibrant on $Arr(\Yy )$.

The conditions (i)$+$(ii)' suffice to prove the desired
result. Indeed, they imply
that in
the diagram $$
\begin{array}{ccc}
\Gamma \underline{Hom}(A,B)&\rightarrow &\underline{Hom}(\Yy ,
nSeCAT')_{1/}(A,B)\\
\downarrow & & \downarrow \\
\Gamma \underline{Hom}(A',B)&\rightarrow &\underline{Hom}(\Yy ,
nSeCAT')_{1/}(A',B)
\end{array}
$$
the bottom arrow is an equivalence; but we know that the two
vertical arrows are equivalences and consequently
the same holds for the top
arrow.


To prove (i) let us fix $\alpha : y\rightarrow z$ in $Arr(\Yy )$
and denote by
$\varphi /\alpha $ the category of the $(y_0,\ldots , z_q)$ in $\Uu $
equipped with a
morphism
from $(y_p\rightarrow z_0)$ to $\alpha$ in $Arr(\Yy )$.
By definition one has
$$
\varphi _{\ast}(F)(\alpha ) = \lim _{\leftarrow , (y,z)\in \varphi /\alpha }
F(y,z).
$$
Now we shall see that the category $\varphi /\alpha $ is obtained by the same
procedure as
$\Uu $ but with other arguments for $Diag$. Let $\Zz$ be the category which
contains the disjoint union of $\Yy /y$ and $z/\Yy$ and which, besides the
arrows of these
categories, has one and only one arrow from each object of $\Yy
/y$ to
each object of $z/\Yy$. One has a morphism of Segal $1$-precategories
$$
\Yy /y \sqcup ^{\{ y\} } I \sqcup ^{\{ z \} } z/\Yy \rightarrow \Zz
$$
and one may verify that this morphism is a weak equivalence.
In particular one has the cofibration
$$
\Yy /y \sqcup z/\Yy \hookrightarrow \Zz ,
$$
and $A|_{\Yy /y}\sqcup B|_{z/\Zz}$ provides a morphism
$$
\Yy /y \sqcup z/\Yy\rightarrow nSeCAT'.
$$
Now, by an argument
analogous to the one given above for $\Uu$,
$Diag (\Zz ,\Yy /y \sqcup z/\Yy, A|_{\Yy /y}\sqcup B|_{z/\Zz})$
is equal to
$\Gamma (\varphi /\alpha , F|_{\varphi /\alpha }).$
By definition of limits, one has
$$
\lim _{(y,z)\in \varphi /\alpha }
F(y,z) = \Gamma (\varphi /\alpha , F|_{\varphi /\alpha }).
$$
Hence the equality
$$
\varphi _{\ast}(F)(\alpha )=
Diag (\Zz ,\Yy /y \sqcup z/\Yy, A|_{\Yy /y}\sqcup B|_{z/\Zz}).
$$

On the other hand, the weak equivalence
$$
\Yy /y \sqcup ^{\{ y\} } I \sqcup ^{\{ z \} } z/\Yy \rightarrow \Zz
$$
gives us
$$
Diag (\Zz ,\Yy /y \sqcup z/\Yy, A|_{\Yy /y}\sqcup B|_{z/\Zz})
\stackrel{\cong}{\rightarrow}
Diag(I, \{ 0,1\} , A(y)\sqcup B(z)).
$$
This last term is identified with $\overline{hom}(A,B)(\alpha )$ (in a way
compatible with the morphisms of source $\overline{hom}(A,B)(\alpha
)$).
This proves that the morphism
$$
\overline{hom}(A,B)(\alpha )\rightarrow
\varphi _{\ast}(F)(\alpha )
$$
is an equivalence of Segal $n$-categories, that is to say (i).

To prove (ii)', we shall show the statement
\newline
(ii)'' the functor
$$
A\mapsto \overline{hom}(A,B)
$$
transforms the cofibrations of the structure of type
HBKQ into fibrations.
\newline
With this result, for an arbitrary $A$,
one may choose an HBKQ-cofibrant replacement
$A'\rightarrow A$, and $\overline{hom}(A',B)$ will be fibrant, which
gives (ii)'. To show (ii)'' it suffices to consider an elementary
cell, i.e. a cofibration of the form
$$
h(X) \times N \rightarrow h(X) \times N'
$$
where $N\rightarrow N'$ is a cofibration of Segal $n$-precategories, and
where $h(X)$ is the functor represented by $X$, namely
$h(X)(Y)=\{ X\rightarrow Y\}$. One must show that
$$
\overline{hom}( h(X) \times N',B)\rightarrow
\overline{hom}( h(X) \times N,B)
$$
is a fibration of Segal $n$-prestacks over $Arr(\Yy )$.
In other words, one must show that if $U\rightarrow U'$ is a trivial
cofibration
of Segal $n$-prestacks over $Arr(\Yy )$,
and if one has morphisms
$$
U\rightarrow \overline{hom}( h(X) \times N',B)
$$
and
$$
U'\rightarrow \overline{hom}( h(X) \times N,B)
$$
which give the same morphism
$U\rightarrow \overline{hom}( h(X) \times N,B)$, then these morphisms
factor through
one and the same lifting: $U'\rightarrow \overline{hom}( h(X) \times N',B)$.

In general, a morphism
$U\rightarrow \overline{hom}(A,B)$ is the data,
for every object $x\rightarrow y$ of $Arr (\Yy )$, of a morphism
$$
U(x\rightarrow y)\times A(x) \rightarrow B(y),
$$
such that, for every morphism
$(x\rightarrow y)\rightarrow (x'\rightarrow y')$ of $Arr (\Yy
)$ (cf above), the two induced morphisms
$$
U(x'\rightarrow y')\times A(x) \rightarrow B(y),
$$
are equal.

Let us denote by $U(X\rightarrow ?)$ the Segal $n$-prestack which is the
inverse image of $U$ under
the morphism $(X/\Yy )^o \rightarrow Arr (\Yy )$; it is a Segal
$n$-prestack on $(X/\Yy )^o$. It follows from the description of the preceding
paragraph that a morphism
$$
U\rightarrow \overline{hom}( h(X) \times N,B)
$$
is nothing other than a morphism of Segal $n$-prestacks on
$(X/\Yy )^o$
$$
U(X\rightarrow ?)\times N \rightarrow B|_{(X/\Yy )^o}.
$$
Our problem therefore becomes: given two morphisms
$$
U(X\rightarrow ?)\times N' \rightarrow B|_{(X/\Yy )^o}
$$
and
$$
U'(X\rightarrow ?)\times N \rightarrow B|_{(X/\Yy )^o}
$$
which coincide on $U(X\rightarrow ?)\times N$,
do they admit a common lifting
to $U'(X\rightarrow ?)\times N'$?

The fact that $B$ is fibrant on $\Yy ^o$  implies that
$B|_{(X/\Yy )^o}$ is fibrant (thanks
to the equality $(X/\Yy )^o=\Yy ^o /X$ and
to Corollary \ref{fibrestrict}). Given an arbitrary cofibration
$N\rightarrow N'$ and a trivial cofibration $U\rightarrow U'$, the
associated morphism $U(X\rightarrow ?)\rightarrow U'(X\rightarrow ?)$
is also a trivial cofibration and one
obtains a trivial cofibration
$$
U(X\rightarrow ?)\times N'\cup ^{U(X\rightarrow ?)\times N}
U'(X\rightarrow ?)\times N \rightarrow U'(X\rightarrow ?)\times N',
$$
which (since
$B|_{(X/\Yy )^o}$ is fibrant) gives the desired lifting property.
\eop

\subnumero{Analysis of \texorpdfstring{$F(A,B;C')$}{F(A,B;C')} on \texorpdfstring{$\Uu$}{U}}
In the above proof we left aside the
problem of
showing that $F(A,B;C')$ is fibrant on $\Uu$, which we
treat now. For this we use the
notion of \emph{Reedy categories},
for which we shall give (in \S 17) only a very
brief presentation: the reader must therefore refer to the
references (Reedy \cite{Reedy}, Bousfield-Kan \cite{BousfieldKan} p. 274,
Dwyer-Kan \cite{DwyerKan3},
Jardine-Goerss \cite{JardineGoerssBook}, Dwyer-Hirschhorn-Kan \cite{DHK},
Hirschhorn \cite{Hirschhorn} Chapter 16). What is given below is
only a rapid argument and
this section is therefore addressed rather to the intrepid reader who already
knows these references.

The category $\Uu$ is a Reedy category for the degree
defined by
$$
deg (y_0,\ldots , y_p; z_0,\ldots , z_q):= p+q;
$$
the ``direct'' morphisms
are
the ``face'' maps which correspond to inclusions of sequences, and
the ``inverse'' morphisms are the degeneracy morphisms
which correspond to surjections in $\Delta \times \Delta$.
The category
opposite $\Uu^o$ is a Reedy category with the same degree
function: the opposite of a direct morphism is inverse
and vice versa.

In particular, for every cmc $M$ the category $M^{\Uu
^o}$ of presheaves of objects of $M$ over $\Uu$ admits a
structure of
cmc which is called the ``Reedy structure'', see \cite{Hirschhorn}.  The
Reedy category $\Uu$ (this is essentially Example 16.1.10 of
\cite{Hirschhorn}) has the following property, analogous to a property of
$\Delta$ (see \cite{Hirschhorn} Corollary 16.4.7):
every
presheaf of sets is cofibrant for the Reedy structure, and in fact
every injection of presheaves of sets on $\Uu$ is a Reedy cofibration
(we leave the proof to the reader). It follows that the same
property holds for the $\Uu$-diagrams in $M$ if $M$ is a cmc of
presheaves of sets whose cofibrations are the injections.
Such is in particular the case for $M=nSePC$.  \footnote{
One could also take an interest in the cmc $nPC$ of non-Segal
$n$-precategories,
defined in \cite{nCAT}.
In that case the cofibrations are the morphisms which are
injective on the objects of $\Theta ^ n$ of non-maximal length; hence our
remark applies and every morphism of $\Uu$-diagrams in $nPC$ which is
an object-by-object cofibration is a Reedy cofibration.}

The category underlying the Reedy cmc $nSePC^{\Uu}$ is
exactly
$nSePCh(\Uu )$; and the weak equivalences (object by object) are the
same. The property stated in the preceding paragraph
says that the cofibrations are the same; hence (\cite{Quillen}) the fibrations
are the same and in fact \emph{the Reedy structure on $nSePC^{\Uu}$
is identified with the cmc structure of \ref{cmf} on $nSePCh(\Uu )$}.

In particular, to verify that $F=F(A,B; C')$ is fibrant
in $nSePCh(\Uu )$
it suffices to verify that $F$ is fibrant for the Reedy structure.
Let us recall what that means (\cite{Hirschhorn} Definition 16.3.2 (3)).
For an object $(y,z)= (y_0,\ldots , y_p; z_0,\ldots , z_q)$ of $\Uu$,
one manufactures a (gigantic) fibre product whose principal terms
(corresponding to sub-sequences of $(y,z)$ of length $p+q-1$)
are the
$$
C'_{p+q-1/}(A(y_0),\ldots , \widehat{A(y_i)},\ldots , A(y_p),
B(z_0),\ldots , B(z_q))
$$
and
$$
C'_{p+q-1/}(A(y_0),\ldots , \ldots , A(y_p),
B(z_0),\ldots ,\widehat{B(z_j)},\ldots , B(z_q)),
$$
which one multiplies together
over the analogous fibre products for the
sub-sequences of length $p+q-2$ (more precisely, the object in question
is expressed as a limit over the ``matching
category'' of $(y,z)\in \Uu $). Let us denote by $Match((y,z); C')$
this (gigantic) Segal $n$-precategory. One has a natural morphism
$$
C'_{p+q/}(A(y_0),\ldots ,  A(y_p),
B(z_0),\ldots , B(z_q))\rightarrow
Match ((y,z); C').
$$
To say that the object we are interested in, namely
$$
(y,z)\mapsto C'_{p+q/}(A(y_0),\ldots ,  A(y_p),
B(z_0),\ldots , B(z_q)),
$$
is fibrant for the Reedy structure,
amounts exactly to saying that the
preceding morphism is a fibration.

Let us denote by $O$ the set
of objects of $C'$ and by $\Delta O$ the category of sequences of objects in
$O$, i.e.
of the $(x_0,\ldots , x_k)$ with $x_i\in O$ (this is the category of simplices
of the category having $O$
as its set of objects and an isomorphism between each pair of objects).
The map $$
\nu _O(C'): (x_0,\ldots , x_k)\mapsto C'_{k/}(x_0,\ldots , x_k)
$$
defines a Segal $n$-prestack over $\Delta O$ and it is easy to
see that if $C'$ is fibrant, then $\nu _O(C')$ is fibrant.
On the other hand $\Delta O$ is also a Reedy category and (as for $\Uu$)
the Reedy structure coincides with the structure of \ref{cmf} on
$nSePCh(\Delta O)$. Hence $\nu _O(C')$ is fibrant for the Reedy structure.
The object $Match ((y,z); C')$ is essentially the same
as the matching object for $\nu _O(C')$
(at least if $p>0$ and
$q>0$---if $p=0$ or $q=0$ there is a missing face whereas if $p=q=0$ everything
is trivial), hence the fact that $\nu
_O(C')$ is fibrant implies that the morphism
$$
C'_{p+q/}(A(y_0),\ldots ,  A(y_p),
B(z_0),\ldots , B(z_q))\rightarrow
Match ((y,z); C')
$$
is fibrant.

Recall that $F$ is defined by the following cartesian diagram
(where the notation is abbreviated):
$$
\begin{array}{ccc}
F(y,z) & \rightarrow & C'_{p+q/}(Ay,
Bz)\\
\downarrow && \downarrow \\
\ast & \rightarrow & C'_{p/}(Ay)\times
C'_{q/}(Bz).
\end{array}
$$
To show that
$F$ is fibrant we must see that the morphism of Segal $n$-prestacks
over $\Uu$
$$
C'_{p+q/}(Ay, Bz)\rightarrow C'_{p/}(Ay)\times
C'_{q/}(Bz)
$$
is a Reedy fibration; for this it is appropriate
to replace the matching object
$$
Match((y,z);C')
$$
above by the \emph{relative matching object}
(\cite{Hirschhorn} Definition 16.2.22),
which we shall denote (for want of a better notation)
$$
Match((y,z);
C'(Ay,Bz)\rightarrow C'(Ay)\times C'(Bz)).
$$
We refer the reader to
\cite{Hirschhorn} Definition 16.2.22 for the definition, which we
do not write out more explicitly (it would be too long).
We remark only that the relative matching map
$$
C'(Ay,Bz)\rightarrow Match((y,z);
C'(Ay,Bz)\rightarrow C'(Ay)\times C'(Bz))
$$
is identical to the map
$$
C'(Ay,Bz)\rightarrow Match((y,z);C')
$$
for $p>0$ and $q>0$; whereas for $p=0$ or $q=0$, in the relative matching
object one recovers exactly the missing face above and the map
$$
C'(Ay,Bz)\rightarrow Match((y,z);
C'(Ay,Bz)\rightarrow C'(Ay)\times C'(Bz))
$$
is equal to the matching map at $(Ay,Bz)$ for $\nu _O(C')$ on
$\Delta O$. In all cases one obtains that the relative matching map
for $C'(Ay,Bz)\rightarrow C'(Ay)\times C'(Bz)$ is a fibration. This proves
that this morphism is a fibration for
the Reedy structure, hence a fibration for the
structure \ref{cmf} on $nSePCh(\Uu )$, and hence that its fibre $F$ is
a fibrant object.
